\documentclass[10pt]{article}
\usepackage[letterpaper,margin=0.85in]{geometry}
\usepackage{amsmath,amssymb,booktabs}
\usepackage[colorlinks=true,urlcolor=blue]{hyperref}
\input{glyphtounicode}
\pdfgentounicode=1
\pdftrailerid{}
\setlength{\emergencystretch}{2em}
\title{Hadamard matrices of orders 5164 and 8332\\from real weighing matrices}
\author{Jeffery Kline}
\date{November 24, 2026\\[3pt]\small Version 0.1.0, release candidate}
\begin{document}
\maketitle
\begin{abstract}
We construct real Hadamard matrices of orders 5164 and 8332, both listed as open in the 2026 survey by Egan, Kotsireas, London and \'{O} Cath\'ain (arXiv:2610.27175v2). For every prime $q\equiv1\pmod4$, we prove that a real weighing matrix of order $q+3$ and weight $q+1$ gives a Hadamard matrix of order $4(q+2)$. Conjugation by a 45-degree rotation in one block of Kumar's array admits the reflection cells that arise from a real weighing matrix. Two weighing matrices constructed by Gysin and Seberry in 1995 then give the stated orders.
\end{abstract}
\section{The construction}
A real Hadamard matrix of order $h$ has entries $\pm1$ and mutually orthogonal rows. A real weighing matrix $W(N,w)$ has entries $0,\pm1$ and satisfies $WW^T=wI_N$. Kumar identified quaternary conference matrices of orders 646 and 1042 as sufficient for the two Hadamard orders considered here. We extend his array to use any real weighing matrix with the parameters below. \href{https://github.com/vatskishan03/Hadamard/blob/d885022144ba6bf4efa8315e14871a3736947527/FRONTIER.md}{Kumar's frontier table}.

\medskip\noindent\textbf{Theorem.} Let $q\equiv1\pmod4$ be prime. A real $W(q+3,q+1)$ gives a real Hadamard matrix of order $4(q+2)$.

\begin{center}
\begin{tabular}{rrrr}\toprule
Golay length & Weighing matrix & Prime $q$ & Hadamard order\\\midrule
320 & $W(1292,1290)$ & 1289 & 5164\\
520 & $W(2084,2082)$ & 2081 & 8332\\\bottomrule
\end{tabular}
\end{center}
\section{From weight deficit two to a conference array}

Write $q+3=2c$, so $c$ is even. Let $W=W(2c,2c-2)$. Each row and column has exactly two zeros since the row and column Grams agree. The zero sets of two rows cannot meet in precisely one column, since their common nonzero support would have odd size. If one row vanishes in columns $a,b$, every row vanishing in $a$ has that same zero set. Column $a$ has exactly two zeros, so exactly two rows share the pair. The distinct pairs partition the columns. Independent row and column permutations put the zero cells on the diagonal of a $c\times c$ array of $2\times2$ cells. This is the pairing in the half-weight construction of \href{https://ajc.maths.uq.edu.au/pdf/12/ocr-ajc-v12-p157.pdf}{Gysin and Seberry, Theorem 3}.

Craigen identified the signed-permutation conference matrix obtained by this compression in 1995. The parity argument below gives a short proof of the cell property. \href{https://us.ticmeet.com/assets/archivos/d6f1d9b8-d39f-4888-925a-7eb81c4905cc/Craigen_handouts.pdf#page=11}{Craigen, Theorem 4 and footnote 16}

Let $h_{ij}$ be the product of the four signs in cell $i,j$, for $i\ne j$. Within row pair $i$, orthogonality gives $c-1$ disagreements, so $\prod_{j\ne i}h_{ij}=(-1)^{c-1}$. Now compare each of those two rows with the same row of pair $k\ne i$. In each comparison there are $2c-4$ common nonzeros and $c-2$ disagreements. Multiplying the two sign products cancels the signs from row $k$, giving $\prod_{j\ne i,k}h_{ij}=1$. Consequently $h_{ik}=(-1)^{c-1}=-1$ for every full cell. Each full cell is a Hadamard matrix of order two.

Put $F=\left(\begin{smallmatrix}1&1\\1&-1\end{smallmatrix}\right)$. Multiplying every cell on the right by $F/2$ gives a real matrix $D$ of Gram $(c-1)I_{2c}$ with zero diagonal cells and off-diagonal cells in the group of eight signed permutation matrices of size two. Conversely multiplying such an array on the right by block-diagonal copies of $F$ recovers a weighing matrix. In particular the input equivalence is exact. Quaternary conference matrices give the special case consisting of quarter-turn rotations.

\section{The modified Kumar array}

Write $n=c-1=(q+1)/2$, $v=2n+1$, and use ordinary transpose on real block matrices throughout. Identify a complex scalar $a+bi$ with $\left(\begin{smallmatrix}a&-b\\ b&a\end{smallmatrix}\right)$. Denote the image of $i$ by $R$, and let $A=(I_2+R)/2$. Thus $A$ represents $(1+i)/2$. Scalars, border signs, identities and the real Yamada circulants below act through this representation. Transpose is taken on both cell coordinates and array coordinates.

Independently multiply block rows and columns of $D$ by signed permutations to make its first block row and column equal to $I_2$ outside the corner. If $Z$ is its remaining $n\times n$ array, its block row and column sums vanish and
\[
ZZ^T=Z^TZ=(nI_n-J_n)\otimes I_2 .
\]
Let $P=\left(\begin{smallmatrix}1&-1\\-1&1\end{smallmatrix}\right)$, $J_2$ the all-ones matrix, and $M,N$ Yamada's real symmetric circulants of order $n$, with $M$ zero on the diagonal, $N$ full, and $M^2+N^2=qI_n$. Use the core and borders from \href{https://github.com/vatskishan03/Hadamard/blob/d885022144ba6bf4efa8315e14871a3736947527/complex-hadamard/CONFERENCE.md#3-the-template}{Kumar, Theorem 3}:
\[
 X'=P\otimes(M\otimes A+N\otimes A^T)
       +J_2\otimes I_n\otimes A,\qquad
 Y'=J_2\otimes Z+P\otimes I_n\otimes I_2 .
\]
In the next display each scalar in a border stands for its $2\times2$ matrix, and $e$ is the all-ones row of length $2n$:
\[
 X=\begin{pmatrix}1&e\\-e^T&X'\end{pmatrix},\qquad
 Y=\begin{pmatrix}1&i e\\-i e^T&Y'\end{pmatrix}.
\]
Form $\widetilde Y$ by replacing every cell $z$ of $Z$ by $A^{-1}zA$. This is global orthogonal conjugation of $Y$ by block copies of $S=\sqrt2 A$, since $S$ is a rotation by $45^\circ$ and every other term in $Y$ commutes with it. Conjugation by $S$ preserves the eight signed permutations: rotations are fixed and the four reflection axes are permuted.

Put
\[
 L=\begin{pmatrix}X&Y\\\widetilde Y^T&-X^T\end{pmatrix}.
\]
Each cell of $L$ is a signed permutation. On an off-diagonal spatial position the cells of $X'$ represent $\pm A\pm A^T$, equal to $\pm I_2$ or $\pm R$. The same possibilities occur on the diagonal using the $J_2$ term instead of $M$. In $Y'$ every off-diagonal spatial position uses a cell of $Z$, and every diagonal position uses $\pm I_2$.

For the diagonal Grams, the core sum is
\[
 X'X'^T+Y'Y'^T
   =\big[(q+2)P\otimes I_n+
           J_2\otimes((2n+1)I_n-2J_n)\big]\otimes I_2
   =(2vI_{2n}-2J_{2n})\otimes I_2.
\]
Here $PJ_2=0$, $P^2=2P$, $J_2^2=2J_2$; the same equality holds with transposes reversed. The block sums of $X'$ are $I_2+R$ in every row and column, and those of $Y'$ are zero. Thus the borders cancel the term $-2J_{2n}$, and their off-diagonal Gram entries vanish by the identical scalar border equations in Kumar's proof. This gives both $XX^T+YY^T=2vI$ and $X^TX+Y^TY=2vI$. Since $X$ commutes with block copies of $S$, $Y^TY=2vI-X^TX$ commutes with them as well. Hence $\widetilde Y^T\widetilde Y=Y^TY$.

For the mixed Gram block $X\widetilde Y-YX$, write $E=e\otimes I_2$, $U=e\otimes R$ and $J=J_{2n}$. The row and column sums of $X'$ are $I_2+R$ and those of $Y',\widetilde Y'$ vanish. Using $R^2=-I_2$, block multiplication gives
\[
\begin{array}{c|cc}
 &X\widetilde Y& YX\\ \hline
(1,1)&I_2+EU^T=I_2-2nR&I_2-UE^T=I_2-2nR\\
(1,2)&U+E\widetilde Y'=U&E+UX'=U\\
(2,1)&-E^T+X'U^T=U^T&U^T-Y'E^T=U^T\\
(2,2)&-J\otimes R+X'\widetilde Y'&-J\otimes R+Y'X'
\end{array}
\]
It remains to show $X'\widetilde Y'=Y'X'$. On the span of $(1,1)$ in the first core coordinate, $X'$ acts as $2I_n\otimes A$, and $Y',\widetilde Y'$ act as $2Z,2\widetilde Z$. Equality follows cellwise from $A\widetilde z=zA$. On the span of $(1,-1)$ both $Y'$ and $\widetilde Y'$ act as twice the identity, which commutes with $X'$. The mixed Gram block therefore vanishes. Thus $LL^T=2vI$. Multiplying $L$ on the right by diagonal copies of $F$ turns each signed permutation into a full sign block and doubles the Gram, giving a real Hadamard matrix of order $4v$.

\section{Golay inputs and the two orders}

For a binary Golay pair $X,Y$ of length $g$, write polynomials in $z$, $\mathcal N(U)=U(z)U(z^{-1})$, so $\mathcal N(X)+\mathcal N(Y)=2g$. Gysin and Seberry's Theorem 7 supplies both $W(4g+4,4g+2)$ and $W(4g+12,4g+10)$. Their Corollary 5 explicitly includes $t=323,521$ for $W(4t,4t-2)$. \href{https://ajc.maths.uq.edu.au/pdf/12/ocr-ajc-v12-p157.pdf#page=6}{Paper, p. 162}

One can check the weights directly. At length $g+1$, take the four words $X+z^g,X-z^g,Y,Y$, padding the last two with zeros. Their norms sum to $4g+2$. At length $g+3$, take $X\pm z^g(1+z-z^2)$ and $Y\pm z^g(1+z^2)$. The two tail norms sum to five, so these four norms sum to $4g+10$. Evaluation on the cyclic shift of the respective length in the Goethals–Seidel array gives the stated weighing matrices. Binary Golay pairs exist at $320=2^5\cdot10$ and $520=2\cdot10\cdot26$ by the classical product construction.

\noindent\textbf{Corollary.} A binary Golay pair of length $g$ gives $H(16g+12)$ whenever $4g+1$ is prime, and $H(16g+44)$ whenever $4g+9$ is prime. For $g=320$ the second formula has $q=1289$. For $g=520$ the first has $q=2081$. Trial division through their square roots proves both primes. Among the 120 orders in Kumar's pinned frontier, the classical lengths $2^a10^b26^d$ with nonnegative exponents yield exactly 5164 and 8332 under these two formulas.


\section{Historical scope and verification}
The entry for each odd part in Seberry and Yamada's 1992 table carries exponent three, giving orders $8\cdot1291=10328$ and $8\cdot2083=16664$. Egan, Kotsireas, London and \'{O} Cath\'ain list both odd parts in their table of open orders $4p$. \href{https://documents.uow.edu.au/~jennie/WEBPDF/175_1992.pdf}{Seberry and Yamada, Table A.2}. \href{https://arxiv.org/pdf/2610.27175v2}{Egan et al., Table 5}.

We located no earlier construction of either order. The full text of Lou, Xia and Xia (2019) was inaccessible during the review. Their abstract describes two special cases for T-matrices of prime order $6m+1$. Every prime of this form admits the representation $a^2+3b^2$ mentioned in the abstract. Whether either special case supplies a T-matrix of order 1291 or 2083 remains unresolved. Such a system would establish the corresponding Hadamard order by an earlier route. We therefore claim constructions at two orders listed as open in Table~5 of arXiv:2610.27175v2. Historical priority remains provisional. \href{https://doi.org/10.17654/DM021020163}{Lou, Xia and Xia (2019)}. \href{https://documents.uow.edu.au/~jennie/WEB/WEB17/Xtb32_6mp1.pdf}{Xia et al. (2017), Theorem 2 and Section 3}.

The accompanying C++ verifier checks all required Golay and Yamada sequence correlations for both applications. Two end-to-end examples at order 76 exercise the two Golay constructions. In a separate reflection example, omitting the conjugation yields 64 failed Gram-block checks. Including it satisfies every check. The C++ verifier uses only one-dimensional inputs at the target orders.

The release target \texttt{make verify-full} constructs each matrix, checks that every entry is a sign, and compares the entire row Gram matrix with $hI_h$ for $h=5164$ and $h=8332$. Each dot product is a sum of $h$ signs. Every binary32 partial sum has integer magnitude at most $h<2^{24}$, so the comparison is exact. The target also checks the SHA-256 digest of each matrix in row-major signed-byte order.

\paragraph{Role of AI.} AI assisted the derivations, the computer verification, the literature search and the exposition under the author's direction. Process-separated AI reviews checked the proofs and certificates; they are not independent expert peer review, and no independent human expert has examined the proofs. The release repository records the sources, the reviews and their dispositions.

\clearpage
\begin{thebibliography}{9}
\bibitem{craigen} R. Craigen, The structure of weighing matrices having large weights,
\emph{Designs, Codes and Cryptography} \textbf{5} (1995), 199--216.
See also \href{https://us.ticmeet.com/assets/archivos/d6f1d9b8-d39f-4888-925a-7eb81c4905cc/Craigen_handouts.pdf}{Signed Groups I}, Theorem 4 and footnote 16.
\bibitem{survey} R. Egan, I. S. Kotsireas, S. London and P. \'{O} Cath\'ain,
\href{https://arxiv.org/abs/2610.27175v2}{Recent progress in the construction of Hadamard matrices}, arXiv:2610.27175v2 (2026), Table 5.
\bibitem{gs} M. Gysin and J. Seberry, \href{https://documents.uow.edu.au/~jennie/WEBPDF/1995_05.pdf}{On the weighing matrices of order $4n$ and weight $4n-2$ and $2n-1$}, \emph{Australasian Journal of Combinatorics} \textbf{12} (1995), 157--174.
\bibitem{kumar} K. Kumar,
\href{https://generalised.si/research/hadamard-matrices}{Hadamard matrices of 34 orders with no known matrix in Table 4 of Cati and Pasechnik}, GSI-R-0001 (2026). The accompanying repository gives the conference construction in Theorem 3.
\bibitem{lxx} L. Lou, M. Xia and T. Xia,
\href{https://doi.org/10.17654/DM021020163}{On T-matrices of order $6m+1$}, \emph{Advances and Applications in Discrete Mathematics} \textbf{21} (2019), 163--181.
\bibitem{sy} J. Seberry and M. Yamada, \href{https://documents.uow.edu.au/~jennie/WEBPDF/175_1992.pdf}{Hadamard matrices, sequences, and block designs}, in J. H. Dinitz and D. R. Stinson (eds.), \emph{Contemporary Design Theory}, Wiley, 1992, 431--560.
\bibitem{yamada} M. Yamada, \href{https://doi.org/10.1017/S144678870003086X}{Some new series of Hadamard matrices}, \emph{Journal of the Australian Mathematical Society, Series A} \textbf{46} (1989), 371--383.
\end{thebibliography}
\end{document}
